\subsection{Extending QuickCheck}
\par Since \ourDSL's internal generator is the same as QuickCheck's, we were able to make it work as an
    extension of QuickCheck. We accomplished this by introducing two new \textit{QuasiQuoters} that
    each parse in their own syntax and translate it into code that can be used in QuickCheck tests.

\par The first new \textit{QuasiQuoter} we introduce is \hskin{[arb| <constructor> <type> = <symbolic expression> |]}.
    During compilation, it reads the string and adds 2 top-level declarations to be used. First, it makes a \hskin{newtype}
    declaration. The name of the new type is the one provided, and it wraps a value of the provided type. Second, it makes
    an instance of QuickCheck's \hskin{Arbitrary} typeclass for this new type. The instance's ``arbitrary'' method generates 
    a value from the provided symbolic expression in the same way that \ourDSL does \evan{this sentence needs to be reworked...}.
    This allows users to use symbolic expressions to tell QuickCheck how to generate values rather than having to write \hskin{Arbitrary}
    instances manually. \evan{probably include an example?}

\begin{figure}
    \begin{smallhskcode}
[arb| SameElem [Int] = x -> x* |]

allSame :: [Int] -> Bool
allSame (x:xs) = all (==x) xs

main :: IO ()
main = quickCheck (\SameElem xs -> allSame xs)
    \end{smallhskcode}
    \caption{\hskin{arb QQ example}}
    \label{fig:arb-qq}
\end{figure}

\begin{figure}
    \begin{smallhskcode}
compose x y = x : y : []
[qc| composeTest :: {x => gt <0>} -> x :: Int -> x :: Int -> x x :: [Int] |]
composeTest (ComposeTest (x, y) f) = x > 0 && x == y & f (compose x y)
    \end{smallhskcode}
    \caption{\hskin{qc QQ example}}
    \label{fig:qc-qq}
\end{figure}

\par The second \textit{QuasiQuoter} we introduce is \hskin{[|qc| <name> :: (<smybolic expression> :: <type> ->)+ <smybolic expression> :: <type> |]}.
    It functions similarly to \hskin{arb} in terms of what declarations are added, but is much more detailed in the exact output.
    This one interprets the name as both the name of a function and a newtype (lowercased and uppercased respectively). It interprets
    each expression-type pair except the last as arguments for the testing function. The last pair is for the output of the testing function.
    The newtype that is declared wraps a singular tuple that contains each of the arguments in order as well as a function that takes the
    provided output type as input and returns a bool.
\par The arbitrary instance for this newtype generates a value for each argument using \ourDSL's generation function and wraps each value in a tuple.
    Then it partially applies arguments to \ourDSL's validation function and passes that as the new type's function argument. This allows
    the user to perform proberty-based tests the way that we designed while using syntax they are more familiar with. It also allows for the
    possibility of adding extra checks before passing the values to the testing function and then the validation functions. \evan{probably also include an example here?}
